\documentclass[12pt,a4paper]{article}
\usepackage[utf8]{inputenc}
\usepackage[T1]{fontenc}
\usepackage[italian]{babel}
\usepackage{libertinus}
\usepackage{amssymb}
\usepackage[a4paper,top=25mm,bottom=25mm,left=22mm,right=22mm,headheight=15pt]{geometry}
\usepackage{graphicx}
\usepackage[table]{xcolor}
\usepackage{tabularx,array}
\usepackage{titlesec}
\usepackage{titletoc}
\usepackage{microtype}
\usepackage{fancyhdr}
\usepackage{enumitem}
\usepackage[hidelinks]{hyperref}

\definecolor{bfblue}{HTML}{1B3F9E}
\definecolor{bfrow}{HTML}{F2F5FC}
\definecolor{bfline}{HTML}{C9D1E6}
\definecolor{bfhead}{HTML}{2E6BE6}
% Segnaposto ancora da compilare (in rosso)
\newcommand{\ph}[1]{\textcolor{red}{[#1]}}
\newcommand{\gl}[1]{\emph{#1}\textsubscript{\scriptsize G}}
\newcommand{\tightlist}{\setlength{\itemsep}{1pt}\setlength{\parskip}{0pt}}
\setlist{topsep=2pt}
\setlist[itemize]{leftmargin=1.3em,label=\textbullet}
\setlist[enumerate]{leftmargin=1.6em}
\setlength{\parindent}{0pt}
\sloppy
\setlength{\emergencystretch}{2em}
\setlength{\parskip}{5pt}
\renewcommand{\arraystretch}{1.45}
\setcounter{secnumdepth}{3}
\setcounter{tocdepth}{3}

\titleformat{\section}{\LARGE\bfseries\color{bfblue}\raggedright}{\thesection.}{0.3em}{}[{\color{bfblue}\vspace{-2pt}\titlerule[0.9pt]}]
\titleformat{\subsection}{\Large\bfseries\color{bfblue}\raggedright}{\thesubsection.}{0.3em}{}
\titleformat{\subsubsection}{\normalsize\bfseries\raggedright}{\thesubsubsection.}{0.3em}{}
\titlespacing*{\section}{0pt}{14pt}{6pt}
\titlespacing*{\subsection}{0pt}{10pt}{4pt}
\titlespacing*{\subsubsection}{0pt}{10pt}{3pt}

% Indice: numeri con punto, trattini fino al numero di pagina
\contentsmargin{2.2em}
\newcommand{\bfleader}{\titlerule*[0.32em]{\textendash}}
\titlecontents{section}[0pt]{\addvspace{5pt}\large\bfseries}{\thecontentslabel.\ }{}{\hspace{0.4em}\bfleader\contentspage}
\titlecontents{subsection}[1.2em]{\addvspace{0pt}}{\thecontentslabel.\ }{}{\hspace{0.4em}\bfleader\contentspage}
\titlecontents{subsubsection}[2.4em]{\addvspace{0pt}}{\thecontentslabel.\ }{}{\hspace{0.4em}\bfleader\contentspage}

\pagestyle{fancy}
\fancyhf{}
\fancyhead[L]{\footnotesize\sffamily\bfseries\color{bfhead}\textls[120]{BENDING FORKS}}
\fancyhead[R]{\footnotesize\sffamily Bill of Materials}
\fancyfoot[C]{\small Pagina \thepage}
\renewcommand{\headrulewidth}{0.9pt}

\newcommand{\bftable}{\rowcolors{2}{bfrow}{white}}
\usepackage{etoolbox}
\AtBeginEnvironment{tabularx}{\bftable}
\usepackage{needspace}
\pretocmd{\section}{\needspace{12\baselineskip}}{}{}
\graphicspath{{./}{../assets/}{../../assets/}{../../../assets/}{assets/}}

% Tabelle lunghe della distinta
\newcolumntype{L}[1]{>{\raggedright\arraybackslash}p{#1}}
\newcolumntype{R}[1]{>{\raggedleft\arraybackslash}p{#1}}
\newcommand{\thc}[1]{\textcolor{white}{\textbf{#1}}}
\newcommand{\rl}{\\ \arrayrulecolor{bfline}\hline}

\begin{document}

% ---------- Copertina ----------
\begin{titlepage}
\centering
\vspace*{4mm}
\includegraphics[width=105mm]{logo-bending-forks-trasparente.png}\\[4mm]
\underline{bending.forks.swe@gmail.com}

\vspace{48mm}
{\fontsize{27}{32}\selectfont\bfseries Bill of Materials\par}
\vspace{3mm}
{\large Distinta dei componenti del dispositivo \emph{FossDot}\par}

\vfill
\begin{tabular}[t]{r|l}
\textbf{Versione} & 0.2.1 \\
\textbf{Uso} & Esterno \\
\textbf{Destinatari} & \emph{Uranion EMEA LTD} \\
 & \emph{Prof. Tullio Vardanega} \\
 & \emph{Prof. Riccardo Cardin} \\
\end{tabular}

\vspace{12mm}
{\large\emph{\textbf{Abstract dei contenuti:}}}\\[1mm]
\begin{minipage}{150mm}\centering Elenco motivato dei componenti da acquistare per il prototipo del dispositivo \emph{FossDot}, con fornitori, costi, tempi di consegna e rischi, come richiesto dall'azienda \emph{Uranion} nella riunione del 13 novembre 2026.\end{minipage}
\vspace*{6mm}
\end{titlepage}

% ---------- Indice ----------
\vspace*{2mm}
{\fontsize{22}{26}\selectfont\bfseries Indice\par}
\vspace{6mm}
\makeatletter\@starttoc{toc}\makeatother
\newpage

% ---------- Contenuto ----------
\section{Registro delle versioni}

\begin{center}
\begin{tabularx}{\linewidth}{>{\hsize=0.55\hsize\raggedright\arraybackslash}X>{\hsize=0.65\hsize\raggedright\arraybackslash}X>{\hsize=0.95\hsize\raggedright\arraybackslash}X>{\hsize=0.95\hsize\raggedright\arraybackslash}X>{\hsize=1.90\hsize\raggedright\arraybackslash}X}
\rowcolor{bfblue}\thc{Versione} & \thc{Data} & \thc{Autore} & \thc{Verificatore} & \thc{Descrizione} \\
0.2.1 & 2026-11-23 & Marco Beltrame & Nicholas Salmaso & Correzioni dalla verifica: verbali con versione nei riferimenti, termini del Glossario segnati nel testo al posto della sezione Glossario, refusi \rl
0.2.0 & 2026-11-18 & Marco Beltrame & Nicholas Salmaso & Scelta della scheda senza \emph{case} e della scocca stampata dal gruppo; prezzi e tempi aggiornati al 18/11; correzioni di forma \rl
0.1.0 & 2026-11-13 & Marco Beltrame, Aldo Toffoli & Nicholas Salmaso & Prima stesura del documento \rl
\end{tabularx}
\end{center}

\section{Introduzione}

\subsection{Scopo del documento}

Il presente documento elenca i componenti da acquistare per realizzare il prototipo del dispositivo \emph{FossDot} e motiva la scelta di ciascuno. Nella riunione del 13 novembre 2026 l'azienda \emph{Uranion} ha chiesto di definire la \emph{Bill of Materials} tra le prime attività del progetto, perché ordinare e ricevere i componenti richiede tempo e senza il dispositivo lo sviluppo del \emph{firmware} diventa difficile.

\subsection{Riferimenti}

\subsubsection{Riferimenti normativi}
\begin{itemize}
\tightlist
\item Norme di progetto v0.1.0;
\item Capitolato d'appalto C3 \emph{FossDot}, \emph{Uranion}: \url{https://www.math.unipd.it/~tullio/IS-1/2026/Progetto/C3-Fossdot.pdf}.
\end{itemize}

\subsubsection{Riferimenti informativi}
\begin{itemize}
\tightlist
\item Verbale esterno del 2026-10-21 con \emph{Uranion}, v1.1 (\emph{wake word} sul dispositivo e linguaggio del \emph{firmware});
\item Verbale interno n.\ 7 del 2026-11-12, v1.0 (requisiti lato \emph{hardware});
\item Verbale interno n.\ 8 del 2026-11-12, v1.0 (incarico di redazione della \emph{Bill of Materials}, VI-8-01);
\item Verbale esterno del 2026-11-13 con \emph{Uranion}, v0.1.0 (richiesta della \emph{Bill of Materials} e numero di microfoni);
\item Glossario v0.2.0;
\item Documentazione dei produttori e pagine dei fornitori elencate nella sezione \ref{sec:fonti}.
\end{itemize}

\section{Vincoli e criteri di scelta}

La scelta dei componenti tiene conto dei requisiti del capitolato e di quanto concordato con l'azienda nei due incontri già svolti. I criteri principali sono i seguenti:

\begin{itemize}
\tightlist
\item \textbf{\gl{Wake word} sul dispositivo}: la parola di attivazione deve essere riconosciuta in locale (verbale del 21/10/2026);
\item \textbf{Numero di microfoni}: due microfoni sono sufficienti per il prodotto, mentre quattro aiutano nella riduzione del rumore e nell'evitare che il dispositivo registri la propria voce, a fronte di un consumo maggiore (DEC-001, verbale del 13/11/2026);
\item \textbf{Costo del dispositivo accettabile}, come richiesto dal capitolato;
\item \textbf{Documentazione completa e reperibile}  (verbale del 13/11/2026);
\item \textbf{Compatibilità verificata} tra processore, microfoni e LED;
\item \textbf{Tempi di consegna brevi}, tenendo conto del rallentamento delle spedizioni nel periodo natalizio;
\item \textbf{\emph{Firmware} preferibilmente in \emph{Rust}}, indicato dal capitolato come valore aggiunto.
\end{itemize}

\subsection{4 microfoni vs 2 microfoni}
\begin{itemize}
\tightlist
\item 4 microfoni (disposizione circolare): copertura reale a 360\textdegree. Permette di calcolare con precisione l'angolo di arrivo della voce (DoA - Direction of Arrival) su tutto il piano orizzontale, ideale se il dispositivo è al centro della stanza (es. su un tavolo o un'isola). Maggiore guadagno di direttività (array gain).
\item 2 microfoni (disposizione lineare): copertura ottimale a 180\textdegree{} frontali. Soffre di ambiguità fronte-retro ed è ideale solo se appoggiato vicino a un muro o su una mensola. Flusso dati dimezzato, minor consumo energetico e gestione molto più semplice.
\end{itemize}

\subsection{C/C++ vs Rust}
Consideriamo valore aggiunto la possibilità di decidere se sviluppare il firmware in C/C++ oppure Rust dato che ad oggi non siamo ancora riusciti a capire quale delle due opzioni sia migliore.


\subsection{DSP vs no DSP}
Il DSP è il responsabile per l'acquisizione dei segnali grezzi dai microfoni, calcola il ritardo di fase (beamforming), abbatte i rumori stazionari (ventole, condizionatore) e esegue l'AEC (Acoustic Echo Cancellation).
In caso di mancanza di DSP la CPU deve campionare N canali in tempo reale, filtrare, allineare le fasi e sottrarre l'eco, aumentando ulteriormente il carico computazionale oltre che complicare la programmazione di essa.
D'altra parte l'assenza di un DSP lascerebbe totale flessibilità nella gestione dell'input audio e un minor consumo complessivo.
Conclusione: riteniamo la presenza del DSP un valore aggiunto poiché semplifica il firmware. In mancanza di un DSP si rischia di rimanere a lungo sulla programmazione del dispositivo a scapito della progressione di altri punti chiave.

\section{Hardware valutato}

Il gruppo ha confrontato le cinque combinazioni proposte nel capitolato e una sesta scheda dello stesso produttore dell'opzione A. La tabella riassume l'esito; il confronto completo è disponibile come documento di supporto.

\begin{center}
\begin{tabularx}{\linewidth}{L{5.2cm}L{2.6cm}X}
\rowcolor{bfblue}\thc{Opzione} & \thc{Esito} & \thc{Motivazione} \\
\emph{reSpeaker XVF3800} + \emph{XIAO ESP32-S3} & Scelta principale & Stesso processore dell'opzione A, elaborazione audio più avanzata e anello di 12 LED e 4 microfoni \rl
A. \emph{reSpeaker Lite} + \emph{XIAO ESP32-S3} & Confronto & Due microfoni come da DEC-001 e costo dimezzato, ma un solo LED e audio meno robusto in ambienti rumorosi \rl
B. \emph{Raspberry Pi 5} + \emph{ReSpeaker USB 4-Mic Array} & Scartata & Costo e consumi elevati, prezzo del \emph{Raspberry Pi 5} alto; la potenza di un sistema \emph{Linux} non è necessaria sul dispositivo \rl
C. \emph{ESP32-S3-Korvo-2} & Scartata & Scheda di sviluppo di grandi dimensioni, cancellazione dell'eco solo in software tramite librerie in C a carico della CPU\rl
D. \emph{Raspberry Pi Pico 2 W} + 4 microfoni MEMS & Scartata & Memoria insufficiente (520 KB) e nessuna elaborazione audio dedicata \rl
E. \emph{nRF5340 Audio DK} + \emph{nRF7002} & Scartata & Costo elevato, \emph{Wi-Fi} su scheda separata \rl
\end{tabularx}
\end{center}

\section{Configurazione proposta}

\subsection{Architettura del dispositivo}

Il dispositivo è composto da una scheda audio con un processore dedicato \emph{XMOS} e da un microcontrollore \emph{ESP32-S3} montato sopra di essa. Il chip \emph{XMOS} riceve il segnale dei microfoni e lo ripulisce da eco, rumore e riverbero; il microcontrollore riceve l'audio già pulito tramite \emph{I2S}, riconosce la \emph{wake word}, invia l'audio al \emph{server} tramite \emph{Wi-Fi} e riproduce la risposta sullo \emph{speaker}. LED e volume sono gestiti dal chip \emph{XMOS}, che il microcontrollore comanda tramite \emph{I2C}.



\subsection{Motivazione della scelta}

Il gruppo propone come scheda principale la \emph{reSpeaker XVF3800} con \emph{XIAO ESP32-S3}. L'azienda ha confermato che due microfoni sono sufficienti, ma si è anche osservato che quattro aiutano a non registrare la voce del dispositivo stesso. Inoltre questa opzione è disponibile in un comodo kit che comprende tutto il materiale necessario per l'assemblaggio del dispositivo, risolvendo così anche la problematica della compatibilità.

Il microcontrollore è lo stesso dell'opzione A: la \emph{wake word} e quasi tutto il \emph{firmware} sono quindi identici per le due schede. Cambiano i seguenti aspetti:
\begin{itemize}
\tightlist
\item Un array di microfoni circolare contro una disposizione lineare;
\item Presenza di XMOS XVF3800 contro un modello precedente e meno performante;
\item DoA, de-reverberation, multi-beamforming in aggiunta alla standard AEC, noise suppression, VNR e AGC;
\item 12 LED invece che 1.
\end{itemize}


\subsection{Caratteristiche della scheda principale}

\begin{center}
\begin{tabularx}{\linewidth}{L{5cm}X}
\rowcolor{bfblue}\thc{Caratteristica} & \thc{Valore} \\
Processore audio & \emph{XMOS XVF3800}: cancellazione dell'eco, \emph{beamforming}, riduzione del rumore, de-riverberazione, controllo automatico del guadagno \rl
Microfoni & 4 microfoni digitali, sensibilità $-26$ dBFS, SNR 64 dBA, \emph{acoustic overload point} 120 dBL \rl
Frequenza di campionamento & 16 kHz massimo, adatta al riconoscimento vocale \rl
Microcontrollore & \emph{XIAO ESP32-S3}: due core a 240 MHz, 8 MB di PSRAM, 8 MB di \emph{flash}, \emph{Wi-Fi} e \emph{Bluetooth} \rl
Uscita audio & connettore per \emph{speaker} fino a 5 W e jack da 3,5 mm \rl
LED & anello di 12 LED RGB indirizzabili \rl
Pulsanti & \emph{mute} dei microfoni in hardware e \emph{reset} \rl
Alimentazione & 5 V da USB o esterna \rl
Dimensioni & 102 $\times$ 102 $\times$ 10 mm \rl
\end{tabularx}
\end{center}

\section{Distinta dei componenti}
I prezzi sono quelli rilevati il 18/11/2026 sul sito \emph{Seeed Studio} (IVA esclusa) e vanno confermati al momento dell'ordine.

\subsection{Gruppo A: dispositivo principale}

\begin{center}
\begin{tabularx}{\linewidth}{L{1.95cm}X L{0.8cm}L{2.5cm}R{1.8cm}R{1.8cm}}
\rowcolor{bfblue}\thc{Codice} & \thc{Componente e modello} & \thc{Q.tà} & \thc{Fornitore} & \thc{Prezzo unit.} & \thc{Totale} \\
BOM-A01 & \emph{reSpeaker XMOS XVF3800} con \emph{XIAO ESP32-S3} (SKU 114993700), senza \emph{case} & 2 & \emph{Seeed Studio} & 60,90 \$ & 121,80 \$ \rl
BOM-A02 & \emph{Speaker} \emph{Seeed Mono Enclosed} 4 $\Omega$ 5 W (SKU 114993346) & 2 & \emph{Seeed Studio} & 2,00 \$ & 4,00 \$ \rl
BOM-A03 & Alimentatore USB-C 5 V 3 A & 2 & {da scegliere} & 10,00 € & 20,00 € \rl
BOM-A04 & Cavo USB-C dati e alimentazione, 1 m & 2 & {da scegliere} & 5,00 € & 10,00 € \rl
BOM-A05 & Scocca progettata dal gruppo e stampata in 3D (circa 100 g di PLA) & 2 & Stampa del gruppo & 2,00 € & 4,00 € \rl
BOM-A06 & Pulsante tattile 6 $\times$ 6 mm con cavetti (opzionale) & 2 & {da scegliere} & 0,50 € & 1,00 € \rl
\end{tabularx}
\end{center}

\begin{center}
\begin{tabularx}{\linewidth}{L{1.95cm}X X}
\rowcolor{bfblue}\thc{Codice} & \thc{Motivazione} & \thc{Alternativa} \\
BOM-A01 & Elaborazione audio dedicata, 4 microfoni con portata fino a 5 m, anello di 12 LED, microcontrollore con memoria sufficiente per la \emph{wake word} locale & \emph{reSpeaker Lite}; versione con \emph{case} dello stesso kit, scartata perché la scocca del gruppo mostra l'anello di LED e contiene lo \emph{speaker} \rl
BOM-A02 & La scheda supporta \emph{speaker} fino a 5 W; stesso produttore della scheda, si ordina insieme alla scheda & \emph{Speaker} generico 4 $\Omega$ 3--5 W, oppure \emph{OpenELAB} (4,19 €) \rl
BOM-A03 & La scheda si alimenta a 5 V; 3 A lasciano margine per \emph{speaker} e LED & Alimentatori USB-C già in possesso del gruppo \rl
BOM-A04 & Programmazione del microcontrollore dalla sua porta USB-C & Cavi già in possesso del gruppo \rl
BOM-A05 & Il capitolato chiede un dispositivo racchiudibile in un contenitore; la scocca del gruppo lascia vedere l'anello di LED e contiene lo \emph{speaker} & Versione con \emph{case} del produttore \rl
BOM-A06 & La scheda non ha pulsanti liberi: serve per ``stop'' o \emph{push-to-talk} & Gestione esclusivamente vocale \rl
\end{tabularx}
\end{center}

\subsection{Gruppo B: dispositivo di confronto}

\begin{center}
\begin{tabularx}{\linewidth}{L{1.95cm}X L{0.8cm}L{2.5cm}R{1.8cm}R{1.8cm}}
\rowcolor{bfblue}\thc{Codice} & \thc{Componente e modello} & \thc{Q.tà} & \thc{Fornitore} & \thc{Prezzo unit.} & \thc{Totale} \\
BOM-B01 & \emph{reSpeaker Lite Voice Assistant Kit} con \emph{XIAO ESP32-S3} saldata (SKU 110061601) & 1 & \emph{Seeed Studio} & 29,90 \$ & 29,90 \$ \rl
BOM-B02 & \emph{Speaker} \emph{Seeed Mono Enclosed} 4 $\Omega$ 5 W & 1 & \emph{Seeed Studio} & 2,00 \$ & 2,00 \$ \rl
BOM-B03 & Alimentatore USB-C 5 V 3 A e cavo dati & 1 & {da scegliere} & 15,00 € & 15,00 € \rl
\end{tabularx}
\end{center}

La \emph{reSpeaker Lite} corrisponde all'opzione A del capitolato: ha due microfoni, come previsto da DEC-001, e costa circa la metà della scheda principale. Serve a verificare con misure concrete se la differenza di qualità audio giustifica il costo.

\section{Riepilogo dei costi}

\begin{center}
\begin{tabularx}{\linewidth}{X R{3.6cm}}
\rowcolor{bfblue}\thc{Voce} & \thc{Solo gruppo A} \\
Schede da \emph{Seeed Studio} & 121,80 \$  \rl
\emph{Speaker} da \emph{Seeed Studio} & 4,00 \$  \rl
Alimentatori, cavi e pulsanti & 31,00 €  \rl
Scocca stampata dal gruppo & 4,00 €  \rl
\textbf{Totale indicativo} & \textbf{125,80 \$ + 35,00 €} \rl
\end{tabularx}
\end{center}

Spedizione, dazi e IVA dell'ordine \emph{Seeed Studio} sono calcolati al momento del pagamento. Il costo di un singolo dispositivo con la scheda principale è di circa 63 \$ (scheda e \emph{speaker}) più 17,50 € di accessori e scocca.

\section{Fornitori e tempi di consegna}

\begin{center}
\begin{tabularx}{\linewidth}{L{2.8cm}L{3.2cm}X L{2cm}}
\rowcolor{bfblue}\thc{Fornitore} & \thc{Componenti} & \thc{Spedizione e tempi} & \thc{Verificato il} \\
\emph{Seeed Studio} & BOM-A01, BOM-A02, BOM-B01, BOM-B02 & Magazzino in Germania per la \emph{XVF3800} (solo UE): al 18/11/2026 su ordinazione, arrivo in magazzino previsto il 20/11/2026, poi spedizione in 1--2 giorni lavorativi. Dal magazzino in Cina 1--3 settimane & 18/11/2026 \rl
\emph{OpenELAB} (Monaco di Baviera) & Alternativa per BOM-A02 e BOM-B02 & Pre-ordine: 10--20 giorni più 3--5 giorni di spedizione & 13/11/2026 \rl
\end{tabularx}
\end{center}

Le schede \emph{reSpeaker} risultano esaurite o in pre-ordine presso diversi rivenditori europei; conviene quindi ordinarle direttamente da \emph{Seeed Studio} appena la distinta è approvata, controllando al momento del pagamento il magazzino di partenza.

\section{Rischi e verifiche all'arrivo}

\begin{center}
\begin{tabularx}{\linewidth}{L{4.2cm}L{3.4cm}X}
\rowcolor{bfblue}\thc{Rischio} & \thc{Effetto} & \thc{Contromisura} \\
Scheda esaurita al momento dell'ordine & Ritardo del \gl{Proof of Concept} & Ordinare appena la distinta è approvata; ripiego sulla \emph{reSpeaker Lite} \rl
Ritardi di spedizione nel periodo natalizio & Ritardo del \emph{Proof of Concept} & Preferire il magazzino in Germania; nel frattempo sviluppo in simulazione \rl
Il \emph{driver} I2S in \emph{Rust} non supporta la ricezione del \emph{clock} dal chip \emph{XMOS} & \emph{Firmware} da adattare & Prova con \emph{esp-hal} all'arrivo della scheda; ripiego su \emph{esp-idf-svc} \rl
Segnale \emph{Wi-Fi} debole con la scocca chiusa & Connessione instabile & Provare la scocca prima della scelta definitiva \rl
\end{tabularx}
\end{center}

All'arrivo dei componenti il gruppo caricherà gli esempi ufficiali del produttore per verificare il funzionamento di microfoni, \emph{speaker}, LED e pulsanti, prima di installare il proprio \emph{firmware}.

\section{Richieste all'azienda}

\begin{enumerate}
\tightlist
\item Approvazione della scheda principale e dell'acquisto della \emph{reSpeaker Lite} per il confronto.
\item Conferma del numero di dispositivi (la proposta è di due più uno).
\item Indicazione di chi effettua l'ordine e dell'indirizzo di consegna dei componenti.
\item Approvazione della scocca progettata dal gruppo al posto della versione con \emph{case} del produttore.
\end{enumerate}

\section{Fonti}\label{sec:fonti}

Pagine consultate il 13, il 15 e il 18 novembre 2026.

\begin{itemize}
\tightlist
\item \emph{Seeed Studio}, \emph{reSpeaker XVF3800} con \emph{XIAO ESP32-S3}: \url{https://www.seeedstudio.com/ReSpeaker-XVF3800-4-Mic-Array-With-XIAO-ESP32S3-p-6489.html}
\item \emph{Seeed Studio}, versione con \emph{case}: \url{https://www.seeedstudio.com/ReSpeaker-XVF3800-With-Case-XIAO-ESP32S3-p-6628.html}
\item \emph{Seeed Studio}, \emph{reSpeaker Lite Voice Assistant Kit}: \url{https://www.seeedstudio.com/ReSpeaker-Lite-Voice-Assistant-Kit-p-5929.html}
\item \emph{Seeed Studio Wiki}, \emph{reSpeaker XVF3800} con \emph{XIAO ESP32-S3}: \url{https://wiki.seeedstudio.com/respeaker_xvf3800_xiao_getting_started/}
\item \emph{Seeed Studio Wiki}, \emph{reSpeaker Lite}: \url{https://wiki.seeedstudio.com/reSpeaker_usb_v3/}
\item \emph{Seeed Studio}, \emph{Mono Enclosed Speaker} 4R 5W: \url{https://www.seeedstudio.com/Mono-Enclosed-Speaker-4R-5W-p-5931.html}
\item \emph{OpenELAB}, \emph{Seeed Studio Mono Enclosed Speaker}: \url{https://openelab.io/products/seeed-studio-mono-enclosed-speaker}
\item \emph{Espressif}, annuncio di \emph{esp-hal} 1.0: \url{https://developer.espressif.com/blog/2025/10/esp-hal-1/}
\end{itemize}

\end{document}
